\documentclass[UTF8,11pt]{ctexart}
\usepackage[a4paper,margin=2.5cm]{geometry}
\usepackage{xcolor}
\usepackage{tcolorbox}
\usepackage{enumitem}
\usepackage{hyperref}
\usepackage{fancyhdr}

% Colors
\definecolor{DeepBlue}{HTML}{065A82}
\definecolor{Teal}{HTML}{1C7293}

% Page style
\pagestyle{fancy}
\fancyhf{}
\fancyhead[L]{船舶水动力学仿真系统 (PINN) 汇报讲稿}
\fancyhead[R]{\thepage}
\fancyfoot[C]{\small 软件工程系学生团队}

% Slide box style
\newtcolorbox{slidebox}[1]{
    colback=DeepBlue!5!white,
    colframe=DeepBlue,
    fonttitle=\bfseries,
    title=#1,
    coltitle=white,
    sharp corners,
    boxrule=0.5pt,
    bottom=2mm,
    top=2mm,
    left=2mm,
    right=2mm
}

\title{\bfseries 船舶水动力学仿真系统 (PINN) 汇报讲稿 \\ \Large （扩充口语化版）}
\author{软件工程系学生团队}
\date{2026年5月16日}

\begin{document}

\maketitle

\section*{基本信息}
\begin{itemize}[label=--]
    \item \textbf{预计演讲时长}：9 -- 10 分钟（字数约 1800 字，按正常语速 180--200 字/分钟计算）
    \item \textbf{适用场景}：软件工程/系统设计课程期末重点汇报
\end{itemize}

\hrulefill

\begin{slidebox}{Slide 1: 封面 (建议时长：40秒)}
各位评委老师、各位同学，大家上午好/下午好！

我是今天的主讲人 [你的名字]，代表我们软件工程小组进行今天的期末项目汇报。我们今天汇报的主题是——《基于 PINN 的船舶水动力学仿真系统需求与架构设计》。

在正式开始之前，我想先简单交代一下背景。当前，整个工业界都在谈论数字化转型和 AI 赋能，但是在船舶工程这样高度依赖物理严谨性的领域，AI 到底能做些什么？它能取代传统的计算工具吗？我们小组带着这些疑问，构思并设计了这套基于前沿人工智能算法的工业软件系统。

接下来，我将按照“提出问题、分析需求、系统设计、以及质量验证”的逻辑主线，为大家完整拆解这套系统的架构与巧思。
\end{slidebox}

\begin{slidebox}{Slide 2: 项目背景与行业痛点 (建议时长：1分30秒)}
首先，第一部分，我们来看看为什么要开发这样一个系统，也就是行业痛点在哪里。

在传统的船舶设计流程中，我们要想知道一艘船在水里开得快不快、阻力大不大，主要依赖 CFD，也就是计算流体力学软件。CFD 技术虽然发展了几十年，精度很高，但在实际的工程应用中，工程师们一直被三个大麻烦困扰：

第一，\textbf{前处理极其复杂}。我们要把一艘船的周围水域划分成几千万个微小的网格，这个过程甚至需要高级工程师耗费数周的时间才能手动调出高质量的网格。
第二，\textbf{计算资源门槛极高}。解算这些庞大的网格需要极其昂贵的超算集群，普通电脑根本跑不动。
第三，\textbf{迭代效率太低了}。如果船型设计师只是把船首的线型稍微改动了一点点，对不起，整个网格划分和计算流程通常要全部推倒重来。

为了打破这个僵局，我们在系统中引入了 \textbf{PINN（物理信息神经网络）} 技术。这是一种非常前沿的 AI 算法。它和大家熟知的 ChatGPT 这类“黑盒”模型不一样。PINN 在训练的时候，会强制把 Navier-Stokes 流体控制方程当作“纪律”或者说“约束”加进去。

它最大的颠覆性在于\textbf{它是无网格的}！系统不需要划分复杂的网格就能直接预测流场。这不仅直接干掉了最耗时的前处理环节，而且一旦模型训练好，遇到新的船型微调，它能做到“秒级”出结果。这就为船型的快速迭代、参数扫描提供了极大的便利。
\end{slidebox}

\begin{slidebox}{Slide 3: 系统建设目标与边界 (建议时长：1分15秒)}
既然技术路线确定了，那么我们这套软件系统的建设目标是什么呢？

我们的核心目标，是打造一个\textbf{完全 Web 化的、高可扩展、并且保证实验绝对可复现}的仿真平台。传统工业软件往往是安装在工作站上的笨重客户端，而我们希望用户只要有个浏览器，就能在云端完成配置、算力调度和结果查看。

这个系统的主要服务对象包括水动力科研人员、造船厂的设计师，甚至是我们在座的高校师生。

为了保证项目在系统设计阶段不跑偏，我们明确了系统的\textbf{边界}：
在“包含范围”内，我们做全套的算例管理、边界条件配置、底层的 PhysicsNeMo 训练任务分发，以及最后的三维可视化和报告生成。
但在“不包含范围”里，我们也必须承认，现阶段它\textbf{不具备}复杂的 CAD 建模修改能力，它的定位也不是为了彻底“淘汰”高保真工业 CFD 软件，而是作为 CFD 的“黄金搭档”，提供一种低成本、高效率的\textbf{前期快速评估工具}。
\end{slidebox}

\begin{slidebox}{Slide 4: 总体架构设计 (解耦与分层) (建议时长：1分30秒)}
明确了需求之后，我们进入核心的软件架构设计部分。请大家看大屏幕。

面对这样一个既有传统 Web 交互，又有重度科学计算的复杂系统，我们架构设计的核心思想只有四个字：\textbf{彻底解耦}。

我们把系统从上到下严格分成了四层：
1. \textbf{最上层是表现层}，基于 React。这里不仅有大家常见的表单和表格，更重要的是，它集成了强大的前端渲染引擎，能够在浏览器里直接流畅地展示三维流场云图。
2. \textbf{第二层是 API 与业务层}，由 Python Flask 构建。大家可以把它理解为系统的“交通枢纽”和“安检站”。它只负责处理用户的权限校验、配置保存和任务下发，绝对不参与任何具体的数学计算。
3. \textbf{第三层是任务调度层}。这也是我们架构设计的一个亮点。我们利用 PostgreSQL 数据库作为任务队列，驱动后台的异步 Worker。为什么要这样？因为一次物理训练短则几小时，长则几天。我们绝对不能让用户的浏览器一直挂着等待。前端发送请求后，立马就能拿到一个任务 ID 关掉页面，后台的 Worker 会在有空闲 GPU 时默默地把活干完。
4. \textbf{最底层是计算引擎层}。我们封装了 NVIDIA 的 PhysicsNeMo 框架，它是整个物理计算的硬核基石。
\end{slidebox}

\begin{slidebox}{Slide 5: 核心业务流程闭环 (建议时长：1分钟)}
好的，看完宏观架构，我们再从用户的视角，走一遍核心的业务流程闭环。这个流程设计的初衷，就是把复杂的物理仿真变成流水线作业。

第一步，\textbf{资产准备}。设计师把船体的几何文件上传到平台，系统会自动进行缩放、包围盒计算等预处理。
第二步，\textbf{物理配置}。用户就像填表一样，输入流体的密度、黏度，告诉系统哪里是入水口，哪里是船体壁面。
第三步，\textbf{模型配置}。算法工程师可以在这里挑选神经网络结构（比如 MLP 或者高阶的 SIREN 网络），然后分配物理方程的权重。
第四步，\textbf{异步执行}。点击提交后，任务进入后台队列，系统会自动调度 GPU 资源开始漫长的训练。用户可以在下班前提交，第二天早上直接看结果。
第五步，也是最有价值的一步，\textbf{结果分析}。系统会自动把生涩的数据转换成直观的速度场、压力场云图，并一键生成对比报告。
\end{slidebox}

\begin{slidebox}{Slide 6: 领域模块设计亮点 (建议时长：1分钟)}
在刚才的流程背后，有几个特别值得一提的模块设计亮点。

首先是 \textbf{Flask API 模块的防阻塞设计}。工业仿真最怕系统卡死，我们的 API 设计严格遵循“快进快出”原则，把重计算全部分离，确保 Web 体验的极度流畅。

其次是 \textbf{Worker 调度模块的物理隔离}。我们把后台打工的 Worker 分成了两类：CPU Worker 和 GPU Worker。生成报告、处理图片这种轻量级工作交给 CPU，昂贵的 GPU 只留给神经网络训练，最大化降低了算力成本。

最后是至关重要的 \textbf{引擎适配层}。大家都知道，AI 框架现在更新迭代的速度简直丧心病狂，可能下个月底层代码的写法就变了。为了保护我们辛苦写的业务代码，我们增加了一层 Adapter（适配器）。不管底层用的是哪个版本的 PhysicsNeMo，只要通过适配器转换，我们的业务逻辑就永远不用改，大大提升了系统的可维护性。
\end{slidebox}

\begin{slidebox}{Slide 7: 数据库与实验可追踪性 (建议时长：1分钟)}
讲完架构，我们来聊聊数据。对于科研类工业软件来说，“不可复现”是致命的。为了解决这个问题，我们设计了\textbf{科研级数据追踪机制}。

最核心的就是 \textbf{配置快照（Snapshot）机制}。当用户点击“开始训练”的那一刻，系统会把当前所有的参数、连同随机种子、依赖库版本，像拍照片一样死死地锁定保存在数据库里。这样一来，无论一年后用户怎么修改他的原配置，历史仿真结果的依据永远都在，100\% 保证实验可复现。

在存储架构上，我们采用了“动静分离”。关系型数据库 PostgreSQL 里面只存轻量级的、结构化的元数据。而几十兆的几何模型、上百兆的神经网络权重以及海量的结果文件，统统放进 MinIO 或 S3 这样的对象存储里。这种设计确保了数据库不会因为大文件拖垮性能，系统架构更加轻盈。
\end{slidebox}

\begin{slidebox}{Slide 8: 质量保证与物理验证 (建议时长：1分15秒)}
接下来，是系统设计中非常独特的一环：质量保证。一般的软件只要代码没有 Bug 就算成功了，但仿真软件如果算出来的结果不符合物理规律，那就等于是制造工业垃圾。因此，我们设计了\textbf{双重质量保障体系}。

第一层，是大家熟悉的\textbf{软件工程验证}。我们有严谨的任务状态机（从排队到运行、到成功或失败，绝不允许越级跳转）；我们还做了严格的项目隔离，并对所有科研数据实行“软删除”，防止误操作带来的灾难。

第二层，是专业的\textbf{物理科学验证}。在 AI 训练时，系统不仅仅看传统的 Loss 曲线有没有下降，更会实时切片分析 PDE（偏微分方程）的残差分布。简单来说，系统在时刻监督这个 AI “有没有违反质量守恒和动量守恒定律”。
不仅如此，我们还设计了一套评级制度。一个模型跑完，必须先在标准的基准船型上测试，只有误差符合工程容忍度，系统才会给它打上 A 级认证，允许它被用于实际方案的评估。
\end{slidebox}

\begin{slidebox}{Slide 9: 总结与演进路线 (建议时长：45秒)}
汇报接近尾声，让我们来做一个总结。

本项目设计的仿真系统，核心价值在于：\textbf{我们利用了现代 Web 高并发架构与前沿的 AI 深度学习技术，成功重塑了传统繁重、低效的船舶水动力仿真工作流。} 它让工业设计变得更敏捷、更普惠。

关于未来的演进路线，我们有着清晰的规划：
目前 \textbf{V1 版本}，我们主要攻克并跑通了静水阻力这个最小核心闭环。
未来的 \textbf{V2 版本}，将重点发力多方案的批量对比与自动化报告生成。
到了 \textbf{V3、V4 阶段}，我们将挑战更复杂的规则波响应模型，并引入 K8s 集群，实现面对海量算力需求的弹性调度能力。
\end{slidebox}

\begin{slidebox}{Slide 10: Q\&A (建议时长：20秒)}
以上就是我们小组关于 PINN 船舶水动力学仿真系统需求与架构设计的全部汇报内容。

感谢各位老师和同学们这一个学期以来的指导与陪伴！希望我们的设计能为大家带来一些工业软件开发的启发。

现在的剩余时间，欢迎各位老师和同学批评指正，进入 Q\&A 提问环节。谢谢大家！
\end{slidebox}

\end{document}
